\documentclass{article}

\title{Indíces de acessibilidade entre distritos do município de São Paulo.}
\author{
	Setorial de Tecnologia e Soberania Digital do PSOL \and
	Júlio César Borges de Oliveira Sampaio
}

% \date{Junho de 2026}
\date{\today}

\usepackage{url}

\usepackage{circuitikz}
\usepackage{graphicx}
\usepackage[portuguese]{babel}

\usepackage[T1]{fontenc}
\usepackage[DefaultFeatures={Ligatures=Common}]{plex-otf} %
% \renewcommand*\familydefault{\sfdefault} %% Only if the base font of the document is to be sans serif

\begin{document}

\maketitle
\pagebreak
\tableofcontents
\listoffigures
\listoftables
\pagebreak

\section{Introdução}
\par
Conforme demandas do movimento Universidade Pública Grajaú, foi feito um pequeno estudo comparativo
de todos os indicadores de acessibilidade em relação aos distritos da cidade de São Paulo, a intenção
é estabelecer critérios comparativos entre os distritos de forma a entender quais tem acessos a determinados
serviços e quais estão mais abandonados pelo serviço público.
\par
No total, os indicadores regionais foram cruzadas com as bordas distritais para tentarmos encontrar indicadores que
mostrem a necessidade de uma instituição de ensino superior no Grajaú, eu tentei o máximo encontrar uma
metodologia que viabilizasse essa classificação.
\section{Os dados}
\subsection{Dados de Acessibilidade}
\par
Os dados mais importante, os indicadores de acessibilidade,
são computados e disponibilizados pelo Instituto de Pesquisa Econômica Aplicada (IPEA),
que nos fornece um arquivo geográfico vetorial (.gpkg) com os dados da cidade, é um dataset
de muito boa qualidade e bem informado, mas
também bem difícil de navegar, em geral, ele divide a cidade em partições hexagonais de formato similar e não por distrito, além disso,
uma determinada divisão pode ter mais de um hexágono, pois existem múltiplas modalidades de minuto e de transporte. Em geral, no
arquivo vetorial existem 163482 `features', que são formas vetoriais, ou seja, hexágonos.
\par
Cada hexágono salva múltiplos tipos de dado, no total são cerca de 50 modalidades se eu não me
engano, essas são explicadas no dicionário de variáveis\cite{dict} e na própria publicação do IPEA\cite{pub}.

\subsection{Distritos}
\subsection{Cidades de São Paulo}
\section{Metodologia}
\subsection{Código e Licença}
\subsection{Computação dos dados distritais}
\subsection{Normalização por População}
\subsection{Normalização por Número de Capturas}
\section{Resultados}
\subsection{Gráficos}
\subsection{Tabelas}
\subsection{Mapas}
\section{Conclusões}

\begin{thebibliography}{9}
\bibitem{dict}
Dicionário de Variáveis (2025) \emph{IPEA} disponível em \url{https://ipea.github.io/aopdata/articles/data_dic_pt.html}.
\bibitem{pub}
Acessibilidade Urbana (2025) \emph{IPEA} disponível em \url{https://www.ipea.gov.br/acessooportunidades/}.

\end{thebibliography}
\end{document}
